% ============================================================================
%  Capitolo 19 -- Istruzioni e interruzioni                      (Lezione 19)
% ============================================================================
\chapter{Istruzioni e interruzioni}
\label{ch:instructions}

\epigraphline{Perché la CPU non riesce mai a combinare niente? Perché viene
interrotta continuamente.}{Battuta anonima sull'architettura dei calcolatori}

\begin{chapgoals}
Alla fine di questo capitolo saprete
\begin{itemize}
\item descrivere un'istruzione come un codice operativo e i suoi operandi, e nominare i tipi di operando;
\item decodificare a mano un'istruzione \riscv, riconoscendone il formato e i campi;
\item seguire un'istruzione attraverso il ciclo del processore: prelievo, decodifica, esecuzione, memoria, scrittura del risultato;
\item spiegare che cos'è una pipeline e perché rende più veloce un processore;
\item definire interruzioni, eccezioni e trap, ed elencarne le cause;
\item descrivere che cosa fa \riscv quando avviene una trap e come riprende il programma interrotto.
\end{itemize}
\end{chapgoals}
% Generated by tools/textbook/export_listings.js from tools/decks/mastery.js and tools/concepts/registry.json: do not edit.
\begin{chapstatus}{Stato di apprendimento di questo capitolo}
\masterychip{now}{DA COMPRENDERE OGGI}\enspace formati delle istruzioni; i modi di indirizzamento; il ciclo del processore; interruzioni e trap.\par\smallskip
\masterychip{use}{DA UTILIZZARE OGGI}\enspace il programma completo, gli strumenti di PRG32 (si usa, non si studia, in questo corso).\par\smallskip
\masterychip{later}{DA COMPRENDERE PIÙ AVANTI}\enspace niente.
\end{chapstatus}
\begin{chapstatus}{Percorso del capitolo}
Alla fine del capitolo saprete:
\begin{enumerate}[nosep,leftmargin=2em]
\item decodificare un'istruzione RISC-V nei suoi campi
\item nominare il modo di indirizzamento di ogni operando
\item descrivere il ciclo dell'istruzione e che cosa cambia un'interruzione
\end{enumerate}
Si appoggia su: solo l'installazione del Capitolo~\ref{ch:setup}. Prossimo passo: Capitolo~\ref{ch:video}.
\end{chapstatus}


% ----------------------------------------------------------------------------
\section{Riscaldamento sulla macchina}

\txtfile[caption={Quattro parole in memoria.},label=lst:ins:words]{code/l19/c03.txt}

Il Listato~\ref{lst:ins:words} è la funzione \code{first\_update} del
Capitolo~\ref{ch:firstcontact} così come sta in memoria: quattro numeri, a
quattro indirizzi di distanza. (La \ins{la} che li precede è omessa, e le
codifiche sono quelle non compresse di 32 bit: con l'estensione C
l'assemblatore può scegliere per alcune di esse forme di 16 bit.) Seguono tre
domande. Come fa il processore a sapere che il primo numero è un load e il
secondo un'addizione? Quali passi compie per ciascuno di essi? E mentre li
esegue, come si accorge che è stato premuto un pulsante o che è scaduto un
timer?

% ----------------------------------------------------------------------------
\section{Le istruzioni}

\subsection{Di che cosa è fatta un'istruzione}

Le istruzioni tipiche di un processore sono trasferimenti di dati tra memoria
e registri, operazioni aritmetiche e logiche, istruzioni per il controllo del
flusso, e trasferimenti di dati da e verso l'ingresso/uscita. (In \riscv la
quarta famiglia non esiste in quanto tale, perché i dispositivi sono memoria.)

\begin{definition}{Istruzione}
Un'\term{istruzione} è composta da un \term{codice operativo}, l'opcode, e da
zero o più \term{operandi}. Un operando può essere una costante (un
\emph{immediato}), una locazione di memoria, o un registro; può essere
esplicito o implicito.
\end{definition}

Le istruzioni possono avere zero, uno, due o tre operandi: \ins{add} ne ha
tre, tutti registri; \ins{ret} non ne ha di visibili.

\subsection{I formati delle istruzioni RISC-V}

Ogni istruzione \riscv dell'insieme di base è lunga 32 bit, ed è disposta
secondo uno di sei \term{formati} (Figura~\ref{fig:ins:formats}). I formati
sono stati progettati in modo che lo stesso campo sia sempre nello stesso
posto: l'opcode nei bit da 6 a 0, il registro destinazione \code{rd} nei bit
da 11 a 7, i registri sorgente \code{rs1} e \code{rs2} nei bit da 19 a 15 e da
24 a 20, il campo \code{funct3} nei bit da 14 a 12. I formati B e J
differiscono da S e U solo nel modo in cui sono disposti i bit dell'immediato.

\begin{figure}[htb]
\centering\footnotesize
\begin{bytefield}[bitwidth=0.0275\textwidth,bitheight=6.5mm]{32}
\bitheader[endianness=big]{0,6,7,11,12,14,15,19,20,24,25,31}\\
\begin{rightwordgroup}{R}\bitbox{7}[bgcolor=prgviolet!20]{funct7} & \bitbox{5}[bgcolor=prgaccent!20]{rs2} & \bitbox{5}[bgcolor=prgaccent!20]{rs1} & \bitbox{3}[bgcolor=prgviolet!20]{funct3} & \bitbox{5}[bgcolor=prggreen!20]{rd} & \bitbox{7}[bgcolor=prgsteel!20]{opcode}\end{rightwordgroup}\\[2pt]
\begin{rightwordgroup}{I}\bitbox{12}[bgcolor=prgamber!30]{imm[11:0]} & \bitbox{5}[bgcolor=prgaccent!20]{rs1} & \bitbox{3}[bgcolor=prgviolet!20]{funct3} & \bitbox{5}[bgcolor=prggreen!20]{rd} & \bitbox{7}[bgcolor=prgsteel!20]{opcode}\end{rightwordgroup}\\[2pt]
\begin{rightwordgroup}{S}\bitbox{7}[bgcolor=prgamber!30]{imm[11:5]} & \bitbox{5}[bgcolor=prgaccent!20]{rs2} & \bitbox{5}[bgcolor=prgaccent!20]{rs1} & \bitbox{3}[bgcolor=prgviolet!20]{funct3} & \bitbox{5}[bgcolor=prgamber!30]{imm[4:0]} & \bitbox{7}[bgcolor=prgsteel!20]{opcode}\end{rightwordgroup}\\[2pt]
\begin{rightwordgroup}{B}\bitbox{7}[bgcolor=prgamber!30]{imm[12$|$10:5]} & \bitbox{5}[bgcolor=prgaccent!20]{rs2} & \bitbox{5}[bgcolor=prgaccent!20]{rs1} & \bitbox{3}[bgcolor=prgviolet!20]{funct3} & \bitbox{5}[bgcolor=prgamber!30]{imm[4:1$|$11]} & \bitbox{7}[bgcolor=prgsteel!20]{opcode}\end{rightwordgroup}\\[2pt]
\begin{rightwordgroup}{U}\bitbox{20}[bgcolor=prgamber!30]{imm[31:12]} & \bitbox{5}[bgcolor=prggreen!20]{rd} & \bitbox{7}[bgcolor=prgsteel!20]{opcode}\end{rightwordgroup}\\[2pt]
\begin{rightwordgroup}{J}\bitbox{20}[bgcolor=prgamber!30]{imm[20$|$10:1$|$11$|$19:12]} & \bitbox{5}[bgcolor=prggreen!20]{rd} & \bitbox{7}[bgcolor=prgsteel!20]{opcode}\end{rightwordgroup}
\end{bytefield}
\caption{I sei formati di \riscv. R: da registro a registro (\ins{add},
\ins{sub}, \ins{and}, \ins{sll}, \ldots). I: immediati e load (\ins{addi},
\ins{andi}, \ins{lw}, \ins{jalr}, \ldots). S: store. B: salti condizionati.
U: immediati superiori (\ins{lui}, \ins{auipc}). J: salti (\ins{jal}).}
\label{fig:ins:formats}
\end{figure}

Questa regolarità è il cuore dell'idea RISC. Il processore sa sempre dove
guardare per prima cosa, e può leggere i registri nominati da un'istruzione
mentre sta ancora scoprendo di quale istruzione si tratta.

\subsection{Decodificare a mano}

\begin{prgconcept}[Il metodo]
Scrivete in binario le otto cifre esadecimali; staccate gli ultimi sette bit,
l'opcode; cercate il formato; poi tagliate i campi.
\end{prgconcept}

\begin{example}{\hex{006283B3}}
$\hex{006283B3}=0000\,0000\,0110\,0010\,1000\,0011\,1011\,0011_2$. Gli ultimi
sette bit, \texttt{0110011}, sono l'opcode delle operazioni da registro a
registro: il formato è R.
\begin{center}\small\ttfamily
\begin{tabular}{|c|c|c|c|c|c|}\hline
0000000&00110&00101&000&00111&0110011\\\hline
\multicolumn{1}{c}{\textrm{\scriptsize funct7}}&\multicolumn{1}{c}{\textrm{\scriptsize rs2 = x6 = t1}}&\multicolumn{1}{c}{\textrm{\scriptsize rs1 = x5 = t0}}&\multicolumn{1}{c}{\textrm{\scriptsize funct3}}&\multicolumn{1}{c}{\textrm{\scriptsize rd = x7 = t2}}&\multicolumn{1}{c}{\textrm{\scriptsize opcode}}
\end{tabular}
\end{center}
Con $\code{funct3}=000$ e $\code{funct7}=0000000$ l'operazione è un'addizione
(con $\code{funct7}=0100000$ sarebbe una sottrazione): l'istruzione è
\code{add t2, t0, t1}.
\end{example}

\begin{example}{\hex{0002A303}}
$\hex{0002A303}=0000\,0000\,0000\,0010\,1010\,0011\,0000\,0011_2$. L'opcode
\texttt{0000011} è quello dei load: il formato è I.
\begin{center}\small\ttfamily
\begin{tabular}{|c|c|c|c|c|}\hline
000000000000&00101&010&00110&0000011\\\hline
\multicolumn{1}{c}{\textrm{\scriptsize immediato = 0}}&\multicolumn{1}{c}{\textrm{\scriptsize rs1 = x5 = t0}}&\multicolumn{1}{c}{\textrm{\scriptsize funct3 = parola}}&\multicolumn{1}{c}{\textrm{\scriptsize rd = x6 = t1}}&\multicolumn{1}{c}{\textrm{\scriptsize opcode}}
\end{tabular}
\end{center}
L'istruzione è \code{lw t1, 0(t0)}, la prima del riscaldamento.
\end{example}

I decodificatori del Capitolo~\ref{ch:combinational} fanno la stessa cosa in
una frazione di nanosecondo. Le tabelle necessarie per decodificare qualsiasi
istruzione sono nell'Appendice~\ref{app:riscv}.

\subsection{Istruzioni a lunghezza fissa o variabile}

In una macchina con istruzioni a \emph{lunghezza fissa} tutte le istruzioni
hanno lo stesso numero di bit, l'istruzione successiva è sempre al PC più una
costante, e prelievo e decodifica sono semplici. I processori RISC sono di
questo tipo: \riscv usa 32 bit e, con l'estensione compressa, 16 o 32,
distinti dai due bit più bassi (Capitolo~\ref{ch:combinational}). In una
macchina con istruzioni a \emph{lunghezza variabile} la lunghezza è nota solo
dopo che la prima parte è stata decodificata: i programmi sono compatti, e la
decodifica è complessa. I processori CISC sono di questo tipo: l'MC68000 usa
da una a cinque parole di 16 bit, e Intel x86 da uno a quindici byte. In
entrambi i casi la lunghezza dell'istruzione decide di quanto avanza il PC.

\subsection{Modi di indirizzamento: dove si trova un operando}

\begin{table}[htb]
\centering\small
\caption{Modi di indirizzamento.}\label{tab:ins:modes}
\begin{tabularx}{\textwidth}{@{}lLl@{}}
\toprule
\textbf{Modo} & \textbf{L'operando è} & \textbf{Esempio \riscv}\\
\midrule
immediato & una costante dentro l'istruzione & \code{addi t1, t1, 1}\\
a registro & il contenuto di un registro & \code{add t2, t0, t1}\\
base più scostamento & la parola di memoria a registro + costante & \code{lw t1, 8(sp)}\\
relativo al PC & un indirizzo a una certa distanza dall'istruzione stessa & \code{beq t0, t1, loop}\\
indiretto in memoria & la parola di memoria il cui indirizzo è in memoria & non disponibile: due load\\
\bottomrule
\end{tabularx}
\end{table}

\riscv ha solo i primi quattro modi della Tabella~\ref{tab:ins:modes}. Un
processore CISC come l'MC68000 ne ha una dozzina, compresi i modi indiretti in
memoria e con autoincremento. L'operazione $R:=M[M[A]]$ diventa in \riscv
\code{lw t0, 0(t1)} seguita da \code{lw t0, 0(t0)}. Il prezzo di modi di
indirizzamento ricchi è un'ulteriore fase nel ciclo del processore, in cui si
assemblano gli operandi.

\subsection{Chiamate, ritorni e la convenzione di chiamata}

Le istruzioni per il controllo del flusso meritano una parola in più, perché
ogni funzione di \vermone dipende da esse. \ins{jal} (\emph{jump and link},
formato J) salta a un indirizzo relativo al PC e scrive nel suo registro
destinazione l'indirizzo dell'istruzione seguente; \ins{jalr} (formato I) fa
lo stesso con una destinazione presa da un registro. La pseudoistruzione
\ins{call} è una \ins{jal} (o una \ins{auipc} con una \ins{jalr}, per
destinazioni lontane) la cui destinazione è \reg{ra}; \ins{ret} è
\code{jalr zero, 0(ra)}.

La \term{convenzione di chiamata} è l'accordo che permette a funzioni scritte
separatamente di chiamarsi l'una con l'altra \citep{riscvabi}. Dice che gli
argomenti viaggiano in \reg{a0}--\reg{a7} e i risultati in \reg{a0} e
\reg{a1}; che una funzione può distruggere i registri \reg{t0}--\reg{t6} e
\reg{a0}--\reg{a7}, che sono quindi \emph{salvati dal chiamante}; che deve
restituire i registri \reg{s0}--\reg{s11} e \reg{sp} come li ha trovati, che
sono quindi \emph{salvati dal chiamato}; e che \reg{sp} è sempre multiplo di
16. È per questo che, fin dal Capitolo~\ref{ch:verm1p2}, un valore che deve
sopravvivere a una chiamata è stato tenuto in un registro \reg{s} e salvato
nel frame. La convenzione è enunciata per intero
nell'Appendice~\ref{app:riscv}.

% ----------------------------------------------------------------------------
\section{Il ciclo del processore}

\subsection{Le fasi del ciclo}

Dopo una fase di avvio, il \emph{bootstrap}, un processore ripete per sempre
la stessa sequenza: il suo algoritmo fondamentale. Nella formulazione RISC
classica la sequenza ha cinque fasi (Figura~\ref{fig:ins:cycle}).

\begin{figure}[htb]
\centering
\begin{tikzpicture}[node distance=4mm]
  \node[blkacc,text width=21mm,minimum height=19mm] (a) {\textbf{Prelievo}\\[1pt]\footnotesize legge l'istruzione; il PC avanza};
  \node[blkamb,text width=21mm,minimum height=19mm,right=of a] (b) {\textbf{Decodifica}\\[1pt]\footnotesize la riconosce; legge i registri};
  \node[blkvio,text width=21mm,minimum height=19mm,right=of b] (c) {\textbf{Esecuzione}\\[1pt]\footnotesize la ALU calcola};
  \node[blkgrn,text width=21mm,minimum height=19mm,right=of c] (d) {\textbf{Memoria}\\[1pt]\footnotesize i load leggono, gli store scrivono};
  \node[blkred,text width=21mm,minimum height=19mm,right=of d] (e) {\textbf{Scrittura}\\[1pt]\footnotesize il risultato va in \texttt{rd}};
  \foreach \x/\y in {a/b,b/c,c/d,d/e} \draw[->,thick] (\x) -- (\y);
  \draw[->,thick,rounded corners=3mm] (e.south) -- ++(0,-6mm) -| (a.south);
\end{tikzpicture}
\caption{Il ciclo del processore.}
\label{fig:ins:cycle}
\end{figure}

\begin{description}[style=nextline,leftmargin=1.5em]
\item[Prelievo (fetch).] L'istruzione all'indirizzo contenuto nel PC viene
  letta dalla memoria; il PC avanza.
\item[Decodifica (decode).] L'opcode viene riconosciuto; i registri nominati
  dall'istruzione vengono letti.
\item[Esecuzione (execute).] La ALU calcola: un risultato, un indirizzo o un
  confronto.
\item[Memoria.] I load leggono la memoria e gli store la scrivono; le altre
  istruzioni saltano questa fase.
\item[Scrittura del risultato (write back).] Il risultato viene scritto nel
  registro destinazione.
\end{description}

\subsection{La fase di prelievo, trasferimento per trasferimento}

Nei termini del modello di riferimento del Capitolo~\ref{ch:processor}, il
prelievo è fatto di quattro \term{trasferimenti tra registri} sul bus interno
(Tabella~\ref{tab:ins:fetch}). In RV32 una parola è di quattro byte e gli
indirizzi contano i byte, così che il PC avanza di 4 (di 2 dopo un'istruzione
compressa).

\begin{table}[htb]
\centering\small
\caption{La fase di prelievo.}\label{tab:ins:fetch}
\begin{tabularx}{\textwidth}{@{}clL@{}}
\toprule
\textbf{Passo} & \textbf{Trasferimento tra registri} & \textbf{Significato}\\
\midrule
1 & $MA := PC$ & l'indirizzo dell'istruzione va al registro indirizzo di memoria\\
2 & abil.\ lettura; $MB := M[MA]$ & la memoria mette l'istruzione sul bus dati; viene catturata nel buffer\\
3 & $D := MB$ & l'istruzione è copiata nel registro di decodifica\\
4 & $PC := PC+4$ & il PC punta ora all'istruzione successiva\\
\bottomrule
\end{tabularx}
\end{table}

\subsection{Il resto del ciclo per due istruzioni}

\begin{table}[htb]
\centering\small
\caption{Le fasi restanti per un'addizione e per un load.}\label{tab:ins:rest}
\begin{tabular}{@{}lll@{}}
\toprule
\textbf{Fase} & \textbf{\code{add t2, t0, t1}} & \textbf{\code{lw t1, 0(t0)}}\\
\midrule
decodifica & $O_1:=R[\reg{t0}]$;\quad $O_2:=R[\reg{t1}]$ & $O_1:=R[\reg{t0}]$;\quad $O_2:=0$ (l'immediato)\\
esecuzione & $U:=O_1+O_2$ & $U:=O_1+O_2$ (un indirizzo)\\
memoria & -- & $MA:=U$; abil.\ lettura; $MB:=M[MA]$\\
scrittura & $R[\reg{t2}]:=U$ & $R[\reg{t1}]:=MB$\\
\bottomrule
\end{tabular}
\end{table}

La stessa ALU somma due registri per \ins{add} e forma un indirizzo per
\ins{lw} (Tabella~\ref{tab:ins:rest}): differiscono solo i segnali di
controllo. L'unità di controllo, decodificando l'opcode, decide quali
trasferimenti avvengono. Uno store termina alla fase di memoria; un salto
condizionato termina alla fase di esecuzione, scegliendo il PC successivo.

\subsection{La pipeline}

Cinque fasi sono cinque pezzi di hardware. Mentre un'istruzione viene
eseguita, la successiva può essere decodificata e una terza prelevata: è una
\term{pipeline} (Tabella~\ref{tab:ins:pipe}).

\begin{table}[htb]
\centering\small
\caption{Una pipeline lineare a cinque stadi: IF prelievo dell'istruzione, ID
decodifica, EX esecuzione, MEM memoria, WB scrittura del risultato.}\label{tab:ins:pipe}
\begin{tabular}{@{}l*{8}{>{\centering\arraybackslash}p{8mm}}@{}}
\toprule
\textbf{ciclo di clock} & 1 & 2 & 3 & 4 & 5 & 6 & 7 & 8\\
\midrule
istruzione 1 & \cellcolor{prgaccent!20}IF & \cellcolor{prgamber!25}ID & \cellcolor{prgviolet!18}EX & \cellcolor{prggreen!20}MEM & \cellcolor{prgred!15}WB &&&\\
istruzione 2 && \cellcolor{prgaccent!20}IF & \cellcolor{prgamber!25}ID & \cellcolor{prgviolet!18}EX & \cellcolor{prggreen!20}MEM & \cellcolor{prgred!15}WB &&\\
istruzione 3 &&& \cellcolor{prgaccent!20}IF & \cellcolor{prgamber!25}ID & \cellcolor{prgviolet!18}EX & \cellcolor{prggreen!20}MEM & \cellcolor{prgred!15}WB &\\
istruzione 4 &&&& \cellcolor{prgaccent!20}IF & \cellcolor{prgamber!25}ID & \cellcolor{prgviolet!18}EX & \cellcolor{prggreen!20}MEM & \cellcolor{prgred!15}WB\\
\bottomrule
\end{tabular}
\end{table}

Ogni istruzione impiega ancora cinque cicli, ma a ogni ciclo di clock si
completa un'istruzione: la produttività è cinque volte maggiore. In generale,
$n$ istruzioni su una pipeline di $k$ stadi si completano in $k+n-1$ cicli. I
processori \term{superscalari} mettono due o più pipeline fianco a fianco, e
completano più di un'istruzione per ciclo (Capitolo~\ref{ch:beyond}). Il core
principale dell'ESP32-C6 ha una pipeline a quattro stadi. I salti, e le
dipendenze tra istruzioni vicine, complicano il quadro: si chiamano
\emph{hazard} (criticità), e sono l'argomento di un corso avanzato.

% ----------------------------------------------------------------------------
\section{Le interruzioni}

\subsection{Il problema del ciclo semplice}

Se il processore facesse sempre e solo prelievo, decodifica ed esecuzione,
nascerebbero tre problemi. L'interazione con il mondo esterno attraverso le
unità di I/O sarebbe inefficiente, perché il programma dovrebbe continuare a
chiedere. Non ci sarebbe modo di rimuovere forzatamente un'applicazione che,
per errore, entra in un ciclo infinito: chi toglie il processore a un
programma bloccato in \code{j loop}? E il sistema operativo avrebbe un
controllo limitato sul sistema.

\subsection{Interruzioni, eccezioni, trap}

\begin{definition}{Interruzione}
Le \term{interruzioni} sono eventi eccezionali, di solito asincroni rispetto
all'esecuzione del programma in corso, che richiedono l'attenzione del
processore.
\end{definition}

La specifica \riscv usa tre parole in modo preciso \citep{riscvpriv}.
Un'\term{interruzione} \eng{interrupt} viene dall'esterno del programma, in un
momento qualsiasi: un timer, un dispositivo. Un'\term{eccezione} è causata
dall'istruzione in esecuzione: un opcode illegale, un indirizzo proibito. Una
\term{trap} è il trasferimento del controllo a un gestore che segue l'una o
l'altra. (I processori più vecchi chiamano «trap» gli eventi sincroni e usano
«interruzione» per tutto.)

\subsection{Il vero ciclo del processore}

Il ciclo reale ha quindi un passo in più (Figura~\ref{fig:ins:real}): dopo
ogni istruzione il processore controlla se c'è un'interruzione in attesa e, in
tal caso, passa al gestore invece di prelevare l'istruzione successiva.

\begin{figure}[htb]
\centering
\begin{tikzpicture}
  \node[blk,fill=prgsteel,text=white,minimum width=22mm,minimum height=9mm] (b) at (0,0) {bootstrap};
  \node[blkacc,minimum width=22mm,minimum height=9mm] (f) at (3.4,0) {prelievo};
  \node[blkamb,minimum width=22mm,minimum height=9mm] (x) at (6.8,0) {decodifica ed\\esecuzione};
  \node[blkred,minimum width=22mm,minimum height=10mm] (q) at (6.8,-1.9) {interruzione\\in attesa?};
  \node[blkvio,minimum width=22mm,minimum height=10mm] (t) at (6.8,-3.8) {trap: si va\\al gestore};
  \draw[->,thick] (b) -- (f); \draw[->,thick] (f) -- (x); \draw[->,thick] (x) -- (q);
  \draw[->,thick] (q.west) -| (f.south) node[pos=0.25,above,lbl]{no};
  \draw[->,thick,prgred] (q) -- (t) node[midway,right,lbl,text=prgred]{sì};
  \draw[->,thick,prgviolet] (t.west) -| (f.south);
  \node[lbl,text=prgviolet] at (4.4,-4.1) {il prossimo prelievo è all'indirizzo in \texttt{mtvec}};
\end{tikzpicture}
\caption{Il vero ciclo del processore.}
\label{fig:ins:real}
\end{figure}

\subsection{Possibili cause di interruzione}

\begin{description}[style=nextline,leftmargin=1.5em]
\item[Interruzione periodica.] Un timer, per esempio ogni 10~ms: per
  condividere il processore tra i programmi e per misurare il tempo che
  ciascuno usa.
\item[Interruzione di I/O.] Una periferica comunica il proprio stato, per
  esempio che è pronta a ricevere dati, per sincronizzarsi con il processore.
\item[Errori nel programma.] Un'istruzione che non esiste o è privilegiata, un
  indirizzo proibito: eventi sincroni, le eccezioni. (In \riscv un overflow
  intero \emph{non} solleva un'eccezione.)
\item[Guasti.] Malfunzionamenti del sistema rilevati da sonde dedicate:
  alimentazione, temperatura, errori di memoria.
\item[Reset.] Riporta l'intero sistema a uno stato noto.
\item[Interruzione software.] Generata da un programma che ha bisogno di una
  risorsa condivisa e chiede al sistema operativo di fare da mediatore: in
  \riscv, l'istruzione \ins{ecall}, \emph{environment call}, che è sincrona
  con il programma che la esegue.
\end{description}

\subsection{Gestire un'interruzione}

\begin{enumerate}
\item Salvare il contesto, parte hardware: almeno l'indirizzo da cui
  riprendere.
\item Identificare la causa dell'interruzione.
\item Saltare al punto di ingresso della \term{routine di servizio
  dell'interruzione} (ISR).
\item Salvare il contesto, parte software: i registri che la routine sta per
  usare.
\item Servire l'interruzione.
\item Se l'interruzione lo consente, ripristinare il contesto e riprendere il
  programma sospeso dall'istruzione che segue l'ultima eseguita.
\end{enumerate}

Il \term{contesto} da salvare e ripristinare è il PC, lo stato e i registri.
Il salvataggio può essere fatto in parte o del tutto in hardware o in
software; \riscv fa il minimo in hardware.

\subsection{Chi ha interrotto? Due soluzioni}

Con più dispositivi il processore deve identificare chi ha interrotto. Con
\emph{registri di stato e polling}, ogni dispositivo ha un registro di stato:
un dispositivo che richiede un'interruzione vi pone a 1 un bit e alza una
linea di richiesta comune; la routine di servizio interroga i dispositivi in
ordine di priorità e serve il primo che trova con il bit alto. Lo schema è
semplice da costruire ma richiede tempo, perché interroga anche dispositivi
che non hanno chiesto niente. Con le \term{interruzioni vettorizzate}
l'identità della sorgente è fornita al processore al momento della richiesta,
e il processore salta direttamente alla routine che le corrisponde: un
meccanismo molto efficiente, in cui il numero di bit del vettore limita il
numero di sorgenti diverse.

La \term{latenza di interruzione} è il ritardo tra l'istante in cui la linea
di richiesta sale e l'istante in cui la routine di servizio comincia a fare
lavoro utile; comprende la fine dell'istruzione corrente, il controllo e la
commutazione.

% ----------------------------------------------------------------------------
\section{La soluzione di RISC-V}

\riscv risolve l'intero problema con pochi registri e un'istruzione.

\subsection{I registri del meccanismo delle trap}

\begin{table}[htb]
\centering\small
\caption{I CSR del modo macchina per il meccanismo delle trap.}\label{tab:ins:csr}
\begin{tabularx}{\textwidth}{@{}lP{0.3\textwidth}L@{}}
\toprule
\textbf{CSR} & \textbf{Scritto da} & \textbf{Contenuto}\\
\midrule
\reg{mtvec} & il firmware, una volta, all'avvio & l'indirizzo del gestore delle trap, e se si usa il modo vettorizzato\\
\reg{mstatus} & hardware e firmware & MIE, l'abilitazione globale delle interruzioni, e MPIE, il suo valore precedente\\
\reg{mie} & il firmware & un bit di abilitazione per ogni sorgente: software, timer, esterna\\
\reg{mip} & l'hardware & un bit di richiesta in attesa per ogni sorgente\\
\reg{mepc} & l'hardware, a ogni trap & l'indirizzo dell'istruzione interrotta\\
\reg{mcause} & l'hardware, a ogni trap & un bit che dice interruzione o eccezione, e un codice per la causa\\
\bottomrule
\end{tabularx}
\end{table}

Un'interruzione viene accolta solo se il suo bit è a 1 sia in \reg{mip} sia in
\reg{mie}, e MIE vale 1 (Tabella~\ref{tab:ins:csr}).

\subsection{Che cosa fa l'hardware su una trap}

\begin{enumerate}
\item $\reg{mepc} :=$ l'indirizzo dell'istruzione interrotta.
\item $\reg{mcause} :=$ la causa: il bit più alto vale 1 per un'interruzione e
  0 per un'eccezione, e gli altri bit contengono un codice.
\item In \reg{mstatus}: $\mathrm{MPIE}:=\mathrm{MIE}$, poi
  $\mathrm{MIE}:=0$, così che ulteriori interruzioni sono disabilitate.
\item Il processore passa al modo macchina.
\item $PC :=$ l'indirizzo contenuto in \reg{mtvec}; oppure, nel modo
  vettorizzato e per le interruzioni, quell'indirizzo più 4 volte la causa.
\end{enumerate}

Alla fine del gestore l'istruzione \ins{mret} disfa tutto:
$\mathrm{MIE}:=\mathrm{MPIE}$ e $PC:=\reg{mepc}$.

L'hardware non inserisce niente su uno stack: il salvataggio dei registri
generici è lasciato al gestore. Il «salvataggio hardware del contesto» dello
schema generale è ridotto a due registri e un bit.

\subsection{La causa di una trap}

\begin{table}[htb]
\centering\small
\caption{Alcuni valori standard di \reg{mcause}; i chip ne aggiungono di propri.}\label{tab:ins:cause}
\begin{tabularx}{\textwidth}{@{}lcL@{}}
\toprule
\textbf{Tipo} & \textbf{Codice} & \textbf{Causa}\\
\midrule
interruzione & 3 & interruzione software di macchina\\
interruzione & 7 & interruzione del timer di macchina\\
interruzione & 11 & interruzione esterna di macchina: un dispositivo, attraverso il controller delle interruzioni\\
eccezione & 2 & istruzione illegale\\
eccezione & 3 & breakpoint (\ins{ebreak}): il modo in cui un debugger ferma un programma\\
eccezione & 5, 7 & errore di accesso in load o in store: un indirizzo proibito\\
eccezione & 8, 11 & environment call (\ins{ecall}) dal modo utente, dal modo macchina\\
\bottomrule
\end{tabularx}
\end{table}

Lo stesso codice 3 significa cose diverse per le interruzioni e per le
eccezioni (Tabella~\ref{tab:ins:cause}): le distingue il bit più alto di
\reg{mcause}. Con \reg{mtvec} nel modo vettorizzato, un'interruzione del timer
entra a $\mathit{base}+4\times 7$.

\subsection{Un gestore delle trap, in schema}

\asmfile[caption={Lo schema di un gestore delle trap.},label=lst:ins:handler]{code/l19/c07.S}

Il programma interrotto non deve accorgersi di niente: ogni registro che il
gestore tocca viene salvato e ripristinato (Listato~\ref{lst:ins:handler}).
Letto come numero con segno, \reg{mcause} è negativo per le interruzioni, così
che \ins{bgez} separa le due famiglie. \ins{mret} è l'unica via di ritorno:
riabilita le interruzioni e ricarica il PC in un solo passo. L'istruzione
\code{csrr t0, mcause} è la pseudoistruzione per
\code{csrrs t0, mcause, zero}.

Il listato è uno schema, non il sorgente di \prg: un gestore reale salva tutti
i registri salvati dal chiamante, di solito su uno stack separato raggiunto
attraverso il registro \reg{mscratch}. Il firmware del kit contiene esattamente
una routine di questo tipo, e una cartuccia viene interrotta molte volte per
frame.

\subsection{Priorità e annidamento}

Per gestire le interruzioni \term{annidate} si stabilisce una priorità: mentre
una routine di servizio è in esecuzione, si accettano solo richieste da
dispositivi di priorità più alta, e una richiesta da un dispositivo ad alta
priorità, come il clock, deve essere accettata anche mentre se ne sta servendo
una di priorità più bassa. La Tabella~\ref{tab:ins:prio} confronta la soluzione
classica dell'MC68000 con quella di \riscv, dove le priorità appartengono a un
controller delle interruzioni esterno al core (sui chip Espressif vanno da 1 a
15).

\begin{table}[htb]
\centering\small
\caption{Priorità e annidamento in due processori.}\label{tab:ins:prio}
\begin{tabularx}{\textwidth}{@{}P{0.22\textwidth}LL@{}}
\toprule
 & \textbf{Motorola MC68000} & \textbf{\riscv}\\
\midrule
priorità di una richiesta & tre linee codificano un livello di priorità dell'interruzione, da 1 a 7 & il controller delle interruzioni dà a ogni sorgente una priorità\\
priorità del processore & tre bit nel registro di stato: la maschera & un registro di soglia nel controller delle interruzioni\\
una richiesta è accettata se & il suo livello è più alto della maschera & la sua priorità è più alta della soglia, ed è abilitata\\
mentre la si serve & la maschera è alzata a quel livello & le interruzioni sono disabilitate ($\mathrm{MIE}=0$) a meno che il gestore non le riabiliti\\
non mascherabile & il livello 7 & un'interruzione non mascherabile, definita da ogni chip\\
identità della sorgente & un numero di vettore di 8 bit dal dispositivo: 256 vettori & \reg{mcause}, poi il controller dice quale dispositivo\\
ritorno & \code{RTE}, che estrae PC e stato dallo stack & \ins{mret}, che usa \reg{mepc} e \reg{mstatus}\\
\bottomrule
\end{tabularx}
\end{table}

\begin{prgwarn}
A volte si dice che «un'interruzione è un segnale che non può essere
mascherato». È vero il contrario: la maggior parte delle interruzioni può
essere mascherata, e solo poche non lo possono.
\end{prgwarn}

\subsection{Le interruzioni sul kit PRG32}

I tre problemi del ciclo semplice, e la loro soluzione, sono tutti visibili
sul kit.

\begin{itemize}
\item Un timer interrompe il processore cento volte al secondo: il firmware lo
  usa per tenere il tempo e per scandire i frame.
\item Le periferiche del display e del suono interrompono quando un blocco di
  dati è stato inviato e serve il successivo.
\item La radio Wi-Fi interrompe quando arriva un pacchetto: una cartuccia può
  essere caricata mentre un gioco è in esecuzione.
\item Le funzioni \code{init}, \code{update} e \code{draw} di una cartuccia
  sono funzioni ordinarie chiamate dal firmware tra un'interruzione e la
  successiva.
\item Una cartuccia bloccata in un ciclo non ferma le interruzioni. \prgqt
  termina il frame con un errore («instruction budget exhausted»); sulla
  scheda un watchdog segnala lo stallo, e l'utente preme reset.
\end{itemize}

In \prgqt un breakpoint è l'eccezione con codice 3.

% ----------------------------------------------------------------------------
% Generated by tools/textbook/export_listings.js from decks/plan.js and code/l19.S: do not edit.
\section{Il programma di questo capitolo}\label{sec:lab:l19}

I listati di questo capitolo sono parti di un solo programma,
\fname{code/l19.S} nel repository del corso, stampato per intero nel
Listato~\ref{lst:lab:l19}. I test automatici del materiale del corso lo
compilano e lo eseguono sull'esecutore senza interfaccia di \prgqt{}, con
ingressi programmati, e controllano il frame; la
Figura~\ref{fig:lab:l19} è il frame scritto dal test.

\begin{description}[style=nextline,leftmargin=1.5em]
\item[Sullo schermo.] Quattro righe di 32 lampade, il codice macchina di lw t1, 0(t0) / addi t1, t1, 1 / sw t1, 0(t0) / ret, e accanto a ciascuna un quadrato verde: la parola in memoria è quella decodificata a mano a lezione.
\item[Che cosa premere.] Niente: il programma va da solo.
\end{description}

\begin{figure}[htb]
\centering
\includegraphics[width=0.5\textwidth]{lab_l19.png}
\caption{\fname{l19.S} su \prgqt{}.}\label{fig:lab:l19}
\end{figure}

\asmfile[caption={Il programma completo del capitolo, \fname{code/l19.S}.},label=lst:lab:l19]{code/l19/program.S}

Per eseguirlo su \prgqt{} i comandi sono quelli del
Listato~\ref{lst:lab:l19:run}; il Paragrafo~\ref{sec:prg:testing} dà gli
stessi passi per QEMU e per la scheda.

\begin{shell}[caption={Compilare \fname{l19.S} ed eseguirlo su \prgqt{}.},label=lst:lab:l19:run]
# 1. start PRG32-QT on this computer (it listens on port 8080)

# 2. in the PRG32 folder, with ESP-IDF active: build the cartridge
python3 -m prg32 cartridge build ~/arc12/code/l19.S --portable \
    --entry-prefix l19 --name l19 --out l19.prg32

# 3. send it to PRG32-QT and start it
python3 -m prg32 esp32c6 upload-and-run l19.prg32 \
    --url http://localhost:8080 --slot cart0

# 4. look, press, compare with the previous slide
\end{shell}


% ----------------------------------------------------------------------------
% Generated by tools/textbook/export_listings.js from tools/concepts/assessments.json: do not edit.
\begin{chapstatus}{Il filo del capitolo}
Il numero 0x00130313 della Lezione 02 si legge ora campo per campo, e il processore che lo preleva può essere interrotto e ripreso. Prossimo passo: i dispositivi che il processore serve.
\par\medskip\textbf{Biglietto d'uscita} (da due a quattro minuti, da soli, su carta):
\begin{enumerate}[nosep,leftmargin=2em]
\item Nominate il modo di indirizzamento di: l'1 in addi t1, t1, 1; la parola letta da lw t1, 8(sp); la destinazione di beq t0, t1, loop; l'operando t0 di add t2, t0, t1.
\item Arriva un'interruzione mentre il processore sta eseguendo un'istruzione. Che cosa succede a quell'istruzione, quale registro ricorda dove riprendere, e come termina il gestore?
\end{enumerate}
\end{chapstatus}


\section{Verificate di aver capito}

\begin{quiz}{In quali bit di un'istruzione \riscv sta l'opcode?}
\item i sette più significativi
\item i sette meno significativi
\item i bit da 14 a 12
\item dipende dal formato
\end{quiz}

\begin{quiz}{Durante il prelievo di un'istruzione RV32 non compressa, il PC viene incrementato di:}
\item 1
\item 2
\item 4
\item 32
\end{quiz}

\begin{quiz}{Quale registro contiene l'indirizzo da cui riprenderà un programma \riscv interrotto?}
\item \reg{ra}
\item \reg{mtvec}
\item \reg{mepc}
\item \reg{mcause}
\end{quiz}

\begin{quiz}{Quale di questi è un'eccezione e non un'interruzione, nei termini di \riscv?}
\item scade un timer
\item viene premuto un tasto
\item viene prelevata un'istruzione illegale
\item arriva un pacchetto dalla rete
\end{quiz}

\section{Esercizi}

\begin{exercise}[prima su carta, poi su \prg]
\begin{enumerate}[nosep]
\item Decodificate a mano \hex{00A28293} e \hex{40628333}: formato, campi,
  assembly.
\item Codificate \code{sub t0, t1, t2} e \code{sw a0, 8(sp)}.
\item Scrivete i trasferimenti tra registri dell'intero ciclo per
  \code{sw t1, 0(t0)} e per \code{beq t0, t1, target}.
\item Disegnate il diagramma della pipeline per sei istruzioni su una pipeline
  a quattro stadi: quando si completa la sesta?
\item In \prgqt mettete un breakpoint in \code{first\_update}, procedete
  passo-passo, e confrontate il codice macchina mostrato con la vostra
  decodifica del riscaldamento.
\end{enumerate}
\expected{Su carta: \code{addi t0, t0, 10} (immediato 10, \code{rs1} =
\reg{x5}, \code{funct3} = 000, \code{rd} = \reg{x5}, opcode 0010011) e
\code{sub t1, t0, t1} (\code{funct7} = 0100000, \code{rs2} = \reg{x6},
\code{rs1} = \reg{x5}, \code{rd} = \reg{x6}); \hex{407302B3} e
\hex{00A12423}; la sesta istruzione si completa al ciclo 9. In \prgqt: le
quattro parole del riscaldamento, o i loro equivalenti compressi.}
\end{exercise}

\section{Glossario}

\begin{glossarytable}
operation code (opcode) & codice operativo\\
operand & operando\\
addressing mode & modo di indirizzamento\\
fetch / decode / execute & prelievo / decodifica / esecuzione\\
pipeline stage & stadio della pipeline\\
context & contesto\\
interrupt & interruzione\\
exception / trap & eccezione / trap\\
interrupt service routine & routine di servizio dell'interruzione\\
interrupt latency & latenza di interruzione\\
to mask an interrupt & mascherare un'interruzione\\
nesting & annidamento\\
\end{glossarytable}

\section{Conclusioni e letture}

Un'istruzione è un codice operativo e zero, uno o più operandi; in \riscv è
sempre di 32 bit, in sei formati regolari. La lunghezza delle istruzioni,
fissa o variabile, determina come si comporta il PC. Il processore ha un
ciclo: preleva un'istruzione, la decodifica, la esegue e scrive il risultato.
Gli eventi eccezionali rompono il ciclo: il processore salva dov'era, serve
l'evento e riprende. Le interruzioni possono essere abilitate, mascherate e
ordinate per priorità; \riscv lo fa con una manciata di registri e
l'istruzione \ins{mret}.

\paragraph{Letture.} \citet{patterson2020}, paragrafi 2.5, 4.1--4.6 e 4.10;
il manuale dell'insieme di istruzioni \riscv, volume~II, «Privileged
Architecture» \citep{riscvpriv}; \emph{Playing with RISC-V}, capitolo~4, su
\prgqt e il suo debugger \citep{prg32textbook}.
